\section{Stato dell'Arte, Analisi della Concorrenza e Analisi ``Make or Buy''}

\subsection{Stato dell'Arte}

Il settore dell'assistenza domiciliare per anziani dispone oggi di numerose soluzioni digitali, ma spesso ciascuna copre un solo bisogno: promemoria farmacologici, teleassistenza, raccolta parametri o dashboard cliniche. A.L.I.C.E. si colloca in una zona intermedia: non vuole sostituire una cartella clinica professionale, ma integrare in una singola applicazione web le azioni quotidiane piu' rilevanti per il monitoraggio remoto.

Le tecnologie moderne rendono questo approccio sostenibile:

\begin{itemize}
    \item \textbf{Framework full-stack:} Next.js permette di gestire pagine, rendering server-side, Server Components, Server Actions e route API nello stesso progetto.
    \item \textbf{Backend-as-a-Service:} Supabase riduce la complessita' del backend fornendo PostgreSQL, autenticazione, API auto-generate e Row Level Security.
    \item \textbf{Interfacce responsive:} l'applicazione e' utilizzabile da browser desktop e tablet, senza installazione.
    \item \textbf{AI come supporto operativo:} una route server-side puo' inviare dati sintetici a un modello esterno e ricevere suggerimenti testuali per l'operatore.
\end{itemize}

\subsection{Analisi della Concorrenza}

Sono stati considerati due competitor principali, gia' presenti nella versione finale del documento dello scorso appello.

\subsubsection{MyTherapy}

MyTherapy e' un'applicazione focalizzata sull'aderenza terapeutica. Gestisce promemoria per i farmaci, consente il tracciamento di parametri vitali e produce report utilizzabili dal personale medico. Il suo limite, rispetto ad A.L.I.C.E., e' la focalizzazione quasi esclusiva sulla terapia farmacologica: non offre una dashboard costruita attorno a esercizi fisici guidati e umore quotidiano con una vista anziano estremamente semplificata.

\subsubsection{ADiLife}

ADiLife e' un sistema professionale di teleassistenza per strutture sanitarie. Offre robustezza infrastrutturale e hardware dedicato, ma risulta piu' costoso e meno adatto a un progetto universitario basato su web app e strumenti open source. Il focus e' piu' orientato all'emergenza e al monitoraggio clinico professionale che alla routine quotidiana del paziente.

\begin{table}[H]
\centering
\footnotesize
\renewcommand{\arraystretch}{1.3}
\begin{tabular}{|p{3cm}|c|c|c|c|}
\hline
\textbf{Piattaforma} & \textbf{Farmaci} & \textbf{Esercizi} & \textbf{Umore} & \textbf{Dashboard operatore} \\
\hline
MyTherapy & \textcolor{green}{\checkmark} & \textcolor{red}{$\times$} & \textcolor{red}{$\times$} & Parziale \\
\hline
ADiLife & \textcolor{green}{\checkmark} & \textcolor{red}{$\times$} & \textcolor{red}{$\times$} & \textcolor{green}{\checkmark} \\
\hline
A.L.I.C.E. & \textcolor{green}{\checkmark} & \textcolor{green}{\checkmark} & \textcolor{green}{\checkmark} & \textcolor{green}{\checkmark} \\
\hline
\end{tabular}
\caption{Confronto sintetico tra A.L.I.C.E. e competitor}
\label{tab:confronto-competitor}
\end{table}

\subsection{Analisi ``Make or Buy''}

\begin{table}[H]
\centering
\renewcommand{\arraystretch}{1.35}
\begin{tabularx}{\textwidth}{|p{3cm}|X|X|}
\hline
\textbf{Criterio} & \textbf{Make} & \textbf{Buy} \\
\hline
Personalizzazione & Interfaccia progettata per due ruoli: operatore e paziente anziano. & Personalizzazione limitata ai vincoli del prodotto commerciale. \\
\hline
Costi & Bassi in fase iniziale grazie a Next.js, React e Supabase. & Licenze e costi di integrazione piu' alti. \\
\hline
Integrazione & Farmaci, esercizi, umore e dashboard nello stesso modello dati. & Spesso richiede piu' strumenti separati. \\
\hline
Controllo tecnico & Pieno controllo del codice, dello schema e delle regole RLS. & Dipendenza dal vendor. \\
\hline
\end{tabularx}
\caption{Analisi comparativa Make vs Buy}
\label{tab:make-or-buy}
\end{table}

La scelta finale e' \textbf{Make}: lo sviluppo interno consente di mantenere il progetto coerente con gli obiettivi didattici, con l'accessibilita' richiesta e con la reale struttura software implementata.

\newpage
